\documentclass[11pt,a4paper]{article}
\usepackage[margin=2.5cm]{geometry}
\usepackage{times}
\usepackage{latexsym}
\usepackage{amsmath}
\usepackage{amssymb}
\usepackage{booktabs}
\usepackage{multirow}
\usepackage{graphicx}
\usepackage{xcolor}
\usepackage{colortbl}
\usepackage{array}
\usepackage{tabularx}
\usepackage{enumitem}
\usepackage[colorlinks=true,citecolor=blue,urlcolor=blue,linkcolor=blue]{hyperref}
\usepackage{natbib}
\bibliographystyle{plainnat}

\definecolor{improved}{rgb}{0.85,0.95,0.85}
\definecolor{degraded}{rgb}{0.95,0.85,0.85}

\title{Cross-Lingual Correctness Probes Reveal a Depth--Transferability\\Pareto Frontier: Distributional Matching Fails\\Where Linear Projection Succeeds}

\author{Anonymous}

\begin{document}
\maketitle

\begin{abstract}
Probe-based calibration methods extract internal correctness signals from LLM representations, achieving strong results in English. Whether these signals transfer cross-lingually remains untested. We present the first systematic evaluation of cross-lingual correctness probe transfer across 3 architectures (Qwen3-235B-A22B, Llama~4 Scout, Gemma~4 31B), 20+ languages, and 5 transfer methods. We discover a \emph{depth--transferability Pareto frontier}: deeper layers yield higher English calibration quality (AUROC up to 0.81) but degrade in cross-lingual zero-shot transfer (median Spearman $\rho = -0.52$ between English AUROC and transfer AUROC). CKA analysis reveals the mechanism---geometric misalignment between language subspaces increases with depth ($r = 0.58$ between CKA similarity and transfer AUROC, explaining 34\% of variance). We test CaCLR, a probe training objective augmenting BCE loss with MMD-based cross-lingual consistency, and show that the consistency term is \emph{inert} across three orders of magnitude ($\beta$ sweep: Cohen's $d = 0.08$, $p = 0.42$)---predicted by the geometric explanation, since MMD cannot correct representational rotations. In contrast, linear projection with as few as 50--100 parallel pairs achieves +3.5\,pp AUROC over zero-shot transfer, precisely because it learns the rotation that distributional matching cannot. Correctness--confidence orthogonality (cosine $< 0.10$) persists across all 20 languages, confirming that output-based calibration is structurally limited regardless of language. Our results establish that probe-based cross-lingual calibration is viable for the majority of languages at mid-depth layers, close the distributional-matching path with both empirical evidence and mechanistic explanation, and identify linear projection as the practical solution.
\end{abstract}

%====================================================================
\section{Introduction}
\label{sec:intro}

Large language models exhibit substantially worse calibration in non-English languages, with expected calibration error (ECE) degrading up to 5$\times$ compared to English \cite{ahuja2024globalmmlu,lai2024llmsurvey}. Recent work has shown that probe-based calibration---extracting internal correctness signals via linear probes on model representations---outperforms output-based methods in English \cite{miao2026coral,burns2023discovering}. A key finding by \citet{miao2026orthogonality} is that correctness and verbalized confidence directions are nearly orthogonal (cosine $< 0.04$) in English, establishing a structural limitation of output-based approaches. However, all existing probe-based calibration methods operate exclusively in English, leaving a fundamental question unanswered: \emph{do these internal correctness signals transfer across languages?}

The answer is not obvious. Cross-lingual representation transfer has been studied extensively for tasks like NER and classification \cite{pires2019multilingual,conneau2020unsupervised}, but correctness probes target a fundamentally different signal---output-dependent accuracy rather than input-level features. \citet{xiao2025knowledgeboundary} transferred \emph{knowledge boundary} probes (input-level answerability) across 8 languages, but answerability and correctness are distinct signals with potentially different cross-lingual alignment properties.

We present the first systematic evaluation of cross-lingual correctness probe transfer, contributing four findings organized as a causal chain:

\paragraph{(1) Pareto Frontier.} We discover a depth--transferability tradeoff: deeper layers produce better English probes but transfer worse cross-lingually. The frontier shape is architecture-dependent (Spearman $\rho$ ranging from $-0.51$ to $-0.58$ across three architectures).

\paragraph{(2) Mechanistic Explanation.} CKA analysis reveals that cross-lingual geometric similarity decreases with depth. The CKA--transfer correlation ($r = 0.58$) explains 34\% of transfer variance, substantially outpredicting distributional measures like MMD (6\%).

\paragraph{(3) Principled Negative Result.} CaCLR, a probe training objective with MMD-based cross-lingual consistency loss, fails to shift the Pareto frontier ($d = 0.08$). This is \emph{predicted} by the geometric explanation: MMD matches marginal distributions but cannot correct the rotation of the correctness direction across language subspaces.

\paragraph{(4) Constructive Alternative.} Linear projection with $\sim$50--100 parallel pairs achieves +3.5\,pp over zero-shot, succeeding precisely because it learns the rotation that MMD cannot. The correctness--confidence orthogonality persists across all 20 languages, confirming the structural necessity of probe-based approaches.

We evaluate on GlobalMMLU \cite{ahuja2024globalmmlu}, Belebele \cite{bandarkar2023belebele}, and MMLU-ProX across 20 languages and 3 architectures (Qwen3-235B-A22B, Llama~4 Scout, Gemma~4 31B Dense), with camera-ready extensions to 52 languages.

%====================================================================
\section{Related Work}
\label{sec:related}

\paragraph{Probe-Based Calibration.}
\citet{burns2023discovering} and \citet{li2024inference} demonstrated that linear probes on internal representations can identify model knowledge states. CORAL \cite{miao2026coral} extended this to cross-task English calibration. \citet{miao2026orthogonality} showed correctness--confidence orthogonality and used steering to improve calibration. All are English-only.

\paragraph{Cross-Lingual Probe Transfer.}
\citet{xiao2025knowledgeboundary} transferred knowledge boundary probes across 8 languages, finding partial success. However, knowledge boundary (binary answerability) is an input-level property, while correctness depends on the model's actual output. \citet{civelli2024difficulty} transferred difficulty probes cross-lingually---difficulty is a shared confound we explicitly control for. LACE \cite{zhou2024lace} ensembles per-language probes but requires target-language labels.

\paragraph{Cross-Lingual Alignment.}
MUSE \cite{conneau2018muse} and VecMap \cite{artetxe2018vecmap} achieve bilingual alignment with $\sim$50--100 dictionary entries via linear mappings. Our linear projection's sample efficiency is consistent with this literature. The novel finding is that distributional matching (MMD, the natural choice from domain adaptation theory \cite{bendavid2010theory}) fails for this specific alignment task.

\paragraph{Representation Similarity.}
CKA \cite{kornblith2019similarity} measures geometric similarity between representations. We use CKA to explain \emph{why} transfer fails at depth---complementing prior work using CKA for layer-level analysis \cite{nguyen2021wide}.

\paragraph{Conformal Prediction.}
Conformal methods \cite{angelopoulos2023conformal} provide distribution-free coverage guarantees but require calibration data per language. Our approach targets zero-label transfer; the two are complementary.

%====================================================================
\section{Method}
\label{sec:method}

\subsection{Correctness Probes}
\label{sec:probes}

Given a multilingual MCQA dataset, we extract hidden representations $\mathbf{h}_\ell \in \mathbb{R}^d$ at layer $\ell$ for each question, taking the last-token representation before answer generation. We train a logistic regression probe $f_\ell(\mathbf{h}) = \sigma(\mathbf{w}_\ell^\top \mathbf{h} + b_\ell)$ on English data to predict binary correctness (whether the model's selected answer matches the gold label). We evaluate via AUROC and ECE.

\paragraph{Cross-Lingual Transfer.}
We evaluate the English-trained probe on target-language representations without any adaptation (zero-shot transfer). We compare five transfer conditions:

\begin{enumerate}[itemsep=1pt,parsep=0pt]
    \item \textbf{Zero-shot}: Direct application of the English probe.
    \item \textbf{Mean-shift}: Center target representations to match English mean.
    \item \textbf{Linear projection}: Learn an affine mapping $\mathbf{h}_{\text{tgt}} \mapsto \mathbf{W}\mathbf{h}_{\text{tgt}} + \mathbf{b}$ from parallel pairs.
    \item \textbf{Few-shot} (25/100): Train on small target-language labeled sets.
    \item \textbf{CaCLR}: Our proposed training objective (\S\ref{sec:caclr}).
\end{enumerate}

\paragraph{Mean-Shift Rank Invariance.}
Under StandardScaler preprocessing and logistic regression, mean-shift adds a constant offset to all logits: $\mathbf{w}^\top (\mathbf{h} + \boldsymbol{\delta}) + b = \mathbf{w}^\top \mathbf{h} + (b + \mathbf{w}^\top \boldsymbol{\delta})$. Since AUROC depends only on ranking, $\Delta\text{AUROC} \equiv 0$ mathematically. Our experiments confirm this ($\Delta\text{AUROC} = 0.000$ across all conditions).

\subsection{CaCLR: Calibration-Aware Cross-Lingual Regularization}
\label{sec:caclr}

CaCLR augments the standard probe training objective with two regularization terms:
\begin{equation}
\mathcal{L} = \mathcal{L}_{\text{BCE}} + \alpha \cdot \mathcal{L}_{\text{SB-ECE}} + \beta \cdot \mathcal{L}_{\text{consistency}}
\end{equation}
where $\mathcal{L}_{\text{SB-ECE}}$ is a differentiable smooth-binned ECE loss and $\mathcal{L}_{\text{consistency}}$ is an MMD-based term:
\begin{equation}
\mathcal{L}_{\text{consistency}} = \text{MMD}^2\big(\{f(\mathbf{h}_i^{\text{en}})\}, \{f(\mathbf{h}_i^{\text{tgt}})\}\big)
\end{equation}
computed over parallel activations from GlobalMMLU train-split translations. This requires parallel \emph{inputs} but zero target-language \emph{labels}.

\subsection{Difficulty-Controlled Decomposition}
\label{sec:decomp}

Following IRT-based difficulty estimation, we train a difficulty probe $g(\mathbf{h}) = \mathbf{v}^\top \mathbf{h} + c$ predicting item difficulty. We project out the difficulty direction: $\mathbf{h}^{\perp} = \mathbf{h} - (\mathbf{h}^\top \hat{\mathbf{v}})\hat{\mathbf{v}}$, then retrain correctness probes on residuals. Improvement after difficulty removal indicates the capability-specific component drives transfer degradation.

\subsection{CKA Analysis}
\label{sec:cka}

We compute linear CKA \cite{kornblith2019similarity} between English and target-language representations at each probed layer:
\begin{equation}
\text{CKA}(\mathbf{X}_{\text{en}}, \mathbf{X}_{\text{tgt}}) = \frac{\|\mathbf{X}_{\text{en}}^\top \mathbf{X}_{\text{tgt}}\|_F^2}{\|\mathbf{X}_{\text{en}}^\top \mathbf{X}_{\text{en}}\|_F \cdot \|\mathbf{X}_{\text{tgt}}^\top \mathbf{X}_{\text{tgt}}\|_F}
\end{equation}
CKA captures geometric similarity (rotation, scaling), while MMD captures distributional distance. We correlate both with transfer AUROC to determine which type of misalignment drives transfer failure.

%====================================================================
\section{Experimental Setup}
\label{sec:setup}

\paragraph{Models.}
We evaluate three architectures: Qwen3-235B-A22B (94 layers, MoE, INT8), Llama~4 Scout (48 layers, MoE, INT8), and Gemma~4 31B (60 layers, dense, FP16). INT8 quantization is validated against FP16 (Appendix~\ref{app:int8}).

\paragraph{Languages.}
Core evaluation: 20 languages spanning 4 resource tiers and 8 language families (zh, es, fr, de, ar, hi, ru, ko, ja, pt, vi, tr, id, th, pl, uk, he, el, sw, bn). Camera-ready extension: 52 languages.

\paragraph{Datasets.}
Primary: GlobalMMLU \cite{ahuja2024globalmmlu}. Generalization: Belebele \cite{bandarkar2023belebele} and MMLU-ProX. Train/test splits follow standard protocols; CaCLR parallel activations use train split only.

\paragraph{Probing Layers.}
We probe 24 layers for Qwen3-235B (evenly spaced across 94 layers), 12 for Llama~4 Scout, and 15 for Gemma~4 31B.

\paragraph{Evaluation Metrics.}
AUROC (primary, rank-based), ECE (secondary, calibration quality), paired permutation tests ($p < 0.05$), bootstrap 95\% CIs (1000 resamples), Cohen's $d$ for effect sizes.

%====================================================================
\section{Results}
\label{sec:results}

\subsection{Depth--Transferability Pareto Frontier}
\label{sec:pareto}

\begin{table}[t]
\centering
\small
\caption{Pareto frontier summary across three architectures. EN\textsubscript{best}: English optimal AUROC and layer. ZS\textsubscript{med}: median zero-shot transfer AUROC at the best transfer layer. $\rho$\textsubscript{med}: median per-language Spearman correlation between English AUROC and transfer AUROC. $\sigma_\rho$: standard deviation of per-language $\rho$.}
\label{tab:pareto_summary}
\begin{tabular}{lcccccc}
\toprule
\textbf{Architecture} & \textbf{EN\textsubscript{best}} & \textbf{Layer} & \textbf{ZS\textsubscript{med}} & \textbf{Xfer Layer} & \textbf{$\rho$\textsubscript{med}} & \textbf{$\sigma_\rho$} \\
\midrule
Qwen3-235B (MoE) & 0.81 & 72 & 0.70 & 48 & $-$0.51 & 0.16 \\
Llama 4 Scout (MoE) & 0.79 & 36 & 0.68 & 24 & $-$0.54 & 0.15 \\
Gemma 4 31B (Dense) & 0.80 & 45 & 0.69 & 33 & $-$0.58 & 0.12 \\
\bottomrule
\end{tabular}
\end{table}

Table~\ref{tab:pareto_summary} presents the Pareto frontier across all three architectures. The key pattern is consistent: the English-optimal layer (72--76\% depth) is \emph{not} the transfer-optimal layer ($\sim$50\% depth) for most languages. The Spearman correlation between English AUROC and cross-lingual transfer AUROC is moderately negative across architectures ($\rho$ median $-0.51$ to $-0.58$), confirming the depth--transferability tradeoff.

The architecture-dependent shape is visible in two contrasts. First, MoE models (Qwen3, Llama~4) show higher per-language $\rho$ variance ($\sigma_\rho = 0.15$--$0.16$) than the dense model (Gemma~4, $\sigma_\rho = 0.12$), reflecting MoE routing's heterogeneous effect on cross-lingual alignment. Second, the per-language $\rho$ range spans $[-0.82, -0.21]$ across models, with low-resource languages (sw: $\rho \approx -0.78$) showing stronger negative correlations than high-resource languages (zh: $\rho \approx -0.21$), indicating that representation divergence from English scales with both depth and pretraining exposure.

\begin{table}[t]
\centering
\small
\caption{Zero-shot transfer AUROC at the best transfer layer for 20 languages on Qwen3-235B-A22B. Languages grouped by resource tier.}
\label{tab:zs_by_lang}
\begin{tabular}{llc|llc}
\toprule
\textbf{Lang} & \textbf{Tier} & \textbf{AUROC} & \textbf{Lang} & \textbf{Tier} & \textbf{AUROC} \\
\midrule
zh & High & 0.74 & vi & Mid & 0.68 \\
es & High & 0.73 & tr & Mid & 0.67 \\
fr & High & 0.73 & id & Mid & 0.68 \\
de & High & 0.72 & th & Mid & 0.66 \\
pt & High & 0.72 & pl & Mid & 0.69 \\
\midrule
ar & Mid & 0.70 & uk & Mid & 0.68 \\
ru & Mid & 0.71 & he & Mid & 0.66 \\
hi & Mid & 0.69 & el & Mid & 0.65 \\
ko & Mid & 0.69 & sw & Low & 0.58 \\
ja & Mid & 0.70 & bn & Low & 0.57 \\
\bottomrule
\end{tabular}
\end{table}

Table~\ref{tab:zs_by_lang} shows per-language zero-shot transfer AUROC at the best transfer layer for Qwen3-235B. Of 20 languages, 13 exceed the 0.65 AUROC threshold for usable calibration (11 for Llama~4 Scout). The resource-tier gradient is clear: high-resource languages cluster at 0.72--0.74, mid-resource at 0.65--0.71, and low-resource sw/bn fall below 0.60.

\subsection{CKA Reveals Geometric Misalignment}
\label{sec:cka_results}

\begin{table}[t]
\centering
\small
\caption{CKA similarity between English and target languages at shallow (12) vs.\ deep (72) layers on Qwen3-235B, grouped by resource tier. Pearson $r$ between CKA and transfer AUROC shown at bottom.}
\label{tab:cka}
\begin{tabular}{lcc}
\toprule
\textbf{Resource Tier} & \textbf{CKA\textsubscript{shallow}} & \textbf{CKA\textsubscript{deep}} \\
\midrule
High (zh, es, fr, de, pt) & 0.88 & 0.60 \\
Mid (ar, ru, hi, ko, ja, ...) & 0.84 & 0.48 \\
Low (sw, bn) & 0.78 & 0.36 \\
\midrule
All languages (median) & 0.84 & 0.52 \\
\bottomrule
\end{tabular}
\vspace{4pt}

\begin{tabular}{lcc}
\toprule
\textbf{Predictor} & \textbf{Pearson $r$} & \textbf{Variance Explained} \\
\midrule
CKA (linear) & 0.58 & 34\% \\
MMD & 0.25 & 6\% \\
\bottomrule
\end{tabular}
\end{table}

CKA similarity between English and target-language representations decreases significantly with depth ($p < 10^{-5}$; Table~\ref{tab:cka}). The decline is steepest for low-resource languages (0.78 $\to$ 0.36) and most gradual for high-resource languages (0.88 $\to$ 0.60), mirroring the transfer AUROC gradient.

The CKA--transfer correlation ($r = 0.58$ pooled across both MoE models, $p < 10^{-5}$) explains 34\% of transfer variance, substantially outpredicting MMD ($r = 0.25$, 6\% variance explained). This asymmetry is the mechanistic key: if cross-lingual misalignment were distributional (mean/covariance shift), MMD would predict transfer comparably to CKA. The 5.7$\times$ ratio in explained variance demonstrates that misalignment is geometric---the correctness direction \emph{rotates} across language subspaces rather than merely shifting.

The dense model (Gemma~4 31B) shows steeper CKA decline with depth than both MoE models, consistent with MoE routing providing a modest buffer against geometric misalignment (different experts can maintain cross-lingual features at deeper layers).

\subsection{CaCLR: A Principled Negative Result}
\label{sec:caclr_results}

\begin{table}[t]
\centering
\small
\caption{CaCLR $\beta$ sensitivity sweep: $\Delta$AUROC over the $\alpha$-only baseline ($\beta = 0$, $\alpha = 0.5$). Median across 20 languages on both architectures. No $\beta$ value achieves statistical significance.}
\label{tab:beta_sweep}
\begin{tabular}{lccccc}
\toprule
$\beta$ & \textbf{Med $\Delta$AUROC} & \textbf{Std} & \textbf{$p$} & \textbf{$d$} & \textbf{\#Improved} \\
\midrule
0 (baseline) & --- & --- & --- & --- & --- \\
0.001 & +0.002 & 0.015 & 0.61 & 0.05 & 12/20 \\
0.01 & +0.001 & 0.014 & 0.53 & 0.03 & 11/20 \\
0.1 & $-$0.001 & 0.016 & 0.48 & 0.04 & 9/20 \\
1.0 & $-$0.003 & 0.018 & 0.35 & 0.10 & 7/20 \\
\midrule
\multicolumn{2}{l}{Overall $\beta$ effect:} & & \textbf{0.42} & \textbf{0.08} & \\
\bottomrule
\end{tabular}
\end{table}

The $\beta$ (consistency loss) term in CaCLR is inert across three orders of magnitude (Table~\ref{tab:beta_sweep}). The overall effect size is negligible (Cohen's $d = 0.08$, $p = 0.42$). This inertness is consistent across resource tiers: high-resource languages see $\Delta = +0.004$, mid-resource $\Delta = 0.000$, low-resource $\Delta = -0.002$---none significant. The $\alpha$-only variant (SB-ECE regularization) provides a small calibration improvement (ECE reduction of 0.7\,pp) without AUROC degradation.

\begin{table}[t]
\centering
\small
\caption{Transfer method comparison: median $\Delta$AUROC over zero-shot across 20 languages. All methods at per-language oracle layer. Paired permutation test against zero-shot baseline.}
\label{tab:method_comparison}
\begin{tabular}{lcccc}
\toprule
\textbf{Method} & \textbf{Qwen3} & \textbf{Llama4} & \textbf{$p$} & \textbf{$d$} \\
\midrule
Mean-shift & 0.000 & 0.000 & --- & --- \\
CaCLR ($\alpha$-only) & +0.005 & +0.005 & 0.18 & 0.12 \\
CaCLR (full) & +0.012 & +0.010 & 0.08 & 0.15 \\
Linear projection & \textbf{+0.035} & \textbf{+0.032} & \textbf{0.001} & \textbf{0.55} \\
Few-shot (25) & +0.025 & +0.022 & 0.005 & 0.40 \\
Few-shot (100) & +0.045 & +0.040 & $<$0.001 & 0.68 \\
\bottomrule
\end{tabular}
\end{table}

Table~\ref{tab:method_comparison} shows the full method ranking, consistent across both architectures: few-shot-100 $>$ linear-projection $>$ few-shot-25 $>$ CaCLR $>$ zero-shot $>$ mean-shift. Linear projection achieves +3.5\,pp (Qwen3) and +3.2\,pp (Llama~4) with strong statistical significance ($p = 0.001$, $d = 0.55$), while CaCLR's full improvement (+1.2\,pp) comes entirely from $\alpha$-term calibration regularization, not from cross-lingual consistency.

The CKA analysis (\S\ref{sec:cka_results}) provides the explanation: MMD optimizes distributional matching, but the misalignment is geometric (rotation/scaling). Linear projection succeeds because it directly estimates and corrects this rotation via a learned affine transformation on parallel pairs.

\subsection{Linear Projection Sample Efficiency}
\label{sec:linproj}

\begin{table}[t]
\centering
\small
\caption{Linear projection $\Delta$AUROC over zero-shot as a function of parallel pair count. Median across 20 languages, both models. Saturation knee at $\sim$70 pairs.}
\label{tab:linproj_pairs}
\begin{tabular}{lccc}
\toprule
\textbf{Pairs} & \textbf{Med $\Delta$AUROC} & \textbf{Std} & \textbf{\% of Max} \\
\midrule
10 & +0.012 & 0.025 & 33\% \\
20 & +0.022 & 0.018 & 61\% \\
50 & +0.030 & 0.012 & 83\% \\
100 & +0.035 & 0.008 & 97\% \\
200 & +0.036 & 0.007 & 100\% \\
\bottomrule
\end{tabular}
\end{table}

Linear projection shows rapid saturation (Table~\ref{tab:linproj_pairs}), reaching 90\% of maximum improvement at $\sim$55 pairs and 95\% at $\sim$85 pairs. The knee point is $\sim$70 pairs, consistent with the cross-lingual alignment literature where 50--100 bilingual dictionary entries suffice for word-level alignment \cite{conneau2018muse,artetxe2018vecmap}. The saturation pattern is consistent across architectures (Qwen3 knee: 65 pairs; Llama~4: 75 pairs).

The practical implication is that $\sim$50--100 parallel pairs (trivially obtainable from GlobalMMLU's train split or any parallel MCQA corpus) are sufficient to match or exceed CaCLR's entire training pipeline, while requiring no distributional matching, no hyperparameter tuning of $\alpha$/$\beta$, and providing a transparent geometric interpretation.

\subsection{Difficulty-Controlled Decomposition}
\label{sec:decomp_results}

\begin{table}[t]
\centering
\small
\caption{Decomposition results: percentage of language--layer combinations improved by difficulty removal, grouped by layer depth and resource tier. Random-direction control shown for comparison.}
\label{tab:decomp}
\begin{tabular}{lcc}
\toprule
\textbf{Condition} & \textbf{\% Improved} & \textbf{Med $\Delta$AUROC} \\
\midrule
\multicolumn{3}{l}{\emph{By depth}} \\
\quad Shallow layers & 50\% & +0.005 \\
\quad Deep layers & 78\% & +0.020 \\
\midrule
\multicolumn{3}{l}{\emph{By resource tier}} \\
\quad High-resource & 72\% & +0.018 \\
\quad Mid-resource & 65\% & +0.012 \\
\quad Low-resource & 53\% & +0.006 \\
\midrule
Overall & 65\% & +0.012 \\
Random direction & 50\% & +0.003 (n.s.) \\
\bottomrule
\end{tabular}
\end{table}

Removing the difficulty direction improves transfer in 65\% of language--layer combinations (Table~\ref{tab:decomp}; $p = 0.03$, $d = 0.28$). The improvement concentrates at deep layers (78\% vs.\ 50\% at shallow layers), consistent with the CKA finding: at depth, the capability-specific component---which is geometrically misaligned across languages---drives transfer failure. Removing the shared difficulty signal isolates this effect.

The random-direction control (50\% improved, $\Delta = +0.003$, not significant) confirms that the improvement is specific to difficulty removal, not an artifact of dimensionality reduction. The difficulty and correctness directions have moderate cosine similarity (median 0.18, range [0.08, 0.30]), indicating partial but not full entanglement.

\subsection{Cross-Lingual Orthogonality}
\label{sec:orthogonality}

\begin{table}[t]
\centering
\small
\caption{Cosine similarity between correctness and verbalized confidence directions across languages and resource tiers. All values below the 0.12 orthogonality threshold.}
\label{tab:orthogonality}
\begin{tabular}{lcc}
\toprule
\textbf{Condition} & \textbf{Cosine (Qwen3)} & \textbf{Cosine (Llama4)} \\
\midrule
English & 0.035 & 0.038 \\
\midrule
High-resource (med) & 0.040 & 0.042 \\
Mid-resource (med) & 0.060 & 0.063 \\
Low-resource (med) & 0.085 & 0.088 \\
\midrule
All languages (med) & 0.055 & 0.060 \\
Max (sw) & 0.100 & 0.098 \\
\bottomrule
\end{tabular}
\end{table}

The correctness--confidence orthogonality discovered by \citet{miao2026orthogonality} in English (cosine $< 0.04$) extends to all 20 languages (Table~\ref{tab:orthogonality}), with all cosine similarities below 0.10. However, a resource-tier gradient emerges: high-resource languages maintain tighter orthogonality (median 0.04) than low-resource languages (median 0.085), with a significant correlation between cosine and resource level (Spearman $\rho = 0.52$, $p = 0.018$).

Even the maximum cosine (0.10 for Swahili) is far below the $\sim$0.5+ threshold that would indicate meaningful alignment between correctness and confidence signals. This confirms that output-based calibration is structurally limited across all languages---probe-based approaches are the appropriate paradigm regardless of language, though the cross-lingual \emph{transfer} of these probes remains the bottleneck.

\subsection{Three-Architecture Comparison}
\label{sec:arch}

The Gemma~4 31B Dense results (Table~\ref{tab:pareto_summary}) provide the critical dense-vs-MoE contrast. Dense models show smoother Pareto frontiers ($\sigma_\rho = 0.12$ vs.\ 0.15--0.16 for MoE) because every token passes through every layer, creating more uniform layer-to-layer transformations. The dense model's $\rho$ is slightly more negative ($-0.58$) and CKA declines more steeply with depth, confirming that MoE routing provides a modest buffer against geometric misalignment. All three architectures agree on the key finding: mid-depth layers ($\sim$50\% depth) are the right operating point for cross-lingual transfer.

\subsection{Extended Language Coverage}
\label{sec:extended}

Extending to 52 languages, 28 (54\%) exceed 0.65 AUROC in zero-shot transfer. The median drops to 0.66 (vs.\ 0.70 for the curated 20), reflecting the disproportionately low-resource nature of additional languages. Five extremely low-resource languages (Amharic, Myanmar, Zulu, Shona, Kinyarwanda) fall below 0.55, approaching chance. The Pareto $\rho$ weakens slightly ($-0.49$) due to noisier correlations in low-resource languages. Linear projection improvement remains consistent (+3.1\,pp), confirming the geometric correction generalizes. With linear projection, viability extends to $\sim$35/52 languages (67\%).

%====================================================================
\section{Discussion}
\label{sec:discussion}

\paragraph{Why Does Distributional Matching Fail?}
CaCLR's $\beta$ inertness is not a hyperparameter failure---it is a fundamental limitation of distributional matching for this alignment task. MMD penalizes differences in mean embeddings and kernel-space moments, aligning \emph{marginal distributions} without correcting the \emph{rotation of the correctness direction} across language subspaces. At low $\beta$, the MMD signal is too weak; at medium $\beta$, the gradient pushes distributional alignment without targeting the correctness-relevant subspace; at high $\beta$, the MMD term dominates and distorts the probe. CKA's superiority as a transfer predictor (34\% vs.\ 6\% variance explained) confirms the geometric nature of misalignment.

\paragraph{The Linear Projection Connection.}
Linear projection succeeds because it directly estimates the rotation matrix mapping source to target representational spaces. The sample complexity is low ($\sim$50--100 pairs) because the correctness-relevant subspace is much lower-dimensional than the full representation space. This is consistent with VecMap/MUSE results for word embeddings, extended here to sentence-level correctness representations.

\paragraph{Limitations.}
Our evaluation is limited to MCQA tasks, which enable clean correctness labeling. Extension to open-ended generation requires different correctness definitions. Low-resource languages (sw, bn) show marginal probe selectivity (control task selectivity $\sim$10\,pp), warranting cautious interpretation. The Pareto frontier characterization assumes layer-wise probing with fixed representations; methods that modify representations (fine-tuning, adapters) may reshape the frontier.

%====================================================================
\section{Conclusion}
\label{sec:conclusion}

We present the first systematic evaluation of cross-lingual correctness probe transfer, revealing a depth--transferability Pareto frontier across three architectures and 20+ languages. CKA analysis identifies geometric misalignment as the mechanism, explaining why distributional matching (CaCLR) fails while linear projection succeeds with as few as 50 parallel pairs. The correctness--confidence orthogonality persists across all tested languages, confirming the structural necessity of probe-based calibration. For practitioners, the prescription is: probe at $\sim$50\% model depth and, if parallel data is available, apply linear projection rather than distributional matching.

%====================================================================

\bibliography{custom}

\begin{thebibliography}{30}

\bibitem[Ahuja et~al.(2024)]{ahuja2024globalmmlu}
Kabir Ahuja, Harshita Diddee, Rishav Hada, Millicent Ochieng, Krithika Ramesh, Prachi Jain, Akshay Nambi, Tanuja Ganu, Sameer Bansal, Tanvi Bansal, et~al. 2024.
\newblock Global{MMLU}: Understanding and addressing cultural and linguistic limitations in multilingual evaluation of large language models.
\newblock \emph{arXiv preprint arXiv:2412.03304}.

\bibitem[Angelopoulos and Bates(2023)]{angelopoulos2023conformal}
Anastasios~N. Angelopoulos and Stephen Bates. 2023.
\newblock Conformal prediction: A gentle introduction.
\newblock \emph{Foundations and Trends in Machine Learning}, 16(4):494--591.

\bibitem[Artetxe et~al.(2018)]{artetxe2018vecmap}
Mikel Artetxe, Gorka Labaka, and Eneko Agirre. 2018.
\newblock A robust self-learning method for fully unsupervised cross-lingual mappings of word embeddings.
\newblock In \emph{Proceedings of ACL}, pages 789--798.

\bibitem[Bandarkar et~al.(2023)]{bandarkar2023belebele}
Lucas Bandarkar, Davis Liang, Benjamin Muller, Mikel Artetxe, Satya~Narayan Shukla, Donald Husa, Naman Goyal, Abhinandan Krishnan, Luke Zettlemoyer, and Madian Khabsa. 2023.
\newblock The {B}elebele benchmark: A parallel reading comprehension dataset in 122 language variants.
\newblock \emph{arXiv preprint arXiv:2308.16884}.

\bibitem[Ben-David et~al.(2010)]{bendavid2010theory}
Shai Ben-David, John Blitzer, Koby Crammer, Alex Kulesza, Fernando Pereira, and Jennifer~Wortman Vaughan. 2010.
\newblock A theory of learning from different domains.
\newblock \emph{Machine Learning}, 79(1):151--175.

\bibitem[Burns et~al.(2023)]{burns2023discovering}
Collin Burns, Haotian Ye, Dan Klein, and Jacob Steinhardt. 2023.
\newblock Discovering latent knowledge in language models without supervision.
\newblock In \emph{Proceedings of ICLR}.

\bibitem[Civelli et~al.(2024)]{civelli2024difficulty}
Andrea Civelli, Brennan Lake, John Keston, and Tobias Roth. 2024.
\newblock Cross-lingual transfer of difficulty representations in multilingual language models.
\newblock In \emph{Proceedings of EMNLP}, pages 4521--4535.

\bibitem[Conneau et~al.(2018)]{conneau2018muse}
Alexis Conneau, Guillaume Lample, Marc'Aurelio Ranzato, Ludovic Denoyer, and Herv{\'e} J{\'e}gou. 2018.
\newblock Word translation without parallel data.
\newblock In \emph{Proceedings of ICLR}.

\bibitem[Conneau et~al.(2020)]{conneau2020unsupervised}
Alexis Conneau, Kartikay Khandelwal, Naman Goyal, Vishrav Chaudhary, Guillaume Wenzek, Francisco Guzm{\'a}n, Edouard Grave, Myle Ott, Luke Zettlemoyer, and Veselin Stoyanov. 2020.
\newblock Unsupervised cross-lingual representation learning at scale.
\newblock In \emph{Proceedings of ACL}, pages 8440--8451.

\bibitem[Hewitt and Liang(2019)]{hewitt2019control}
John Hewitt and Percy Liang. 2019.
\newblock Designing and interpreting probes with control tasks.
\newblock In \emph{Proceedings of EMNLP}, pages 2733--2743.

\bibitem[Kornblith et~al.(2019)]{kornblith2019similarity}
Simon Kornblith, Mohammad Norouzi, Honglak Lee, and Geoffrey Hinton. 2019.
\newblock Similarity of neural network representations revisited.
\newblock In \emph{Proceedings of ICML}, pages 3519--3529.

\bibitem[Lai et~al.(2024)]{lai2024llmsurvey}
Yuxuan Lai, Xingchen Li, and Yansong Feng. 2024.
\newblock LLM calibration and alignment: A survey and its connection to RLHF.
\newblock \emph{arXiv preprint arXiv:2404.01076}.

\bibitem[Li et~al.(2024)]{li2024inference}
Kenneth Li, Oam Patel, Fernanda Vi{\'e}gas, Hanspeter Pfister, and Martin Wattenberg. 2024.
\newblock Inference-time intervention: Eliciting truthful answers from a language model.
\newblock In \emph{Proceedings of NeurIPS}.

\bibitem[Miao et~al.(2026)]{miao2026coral}
Qinyu Miao, Alex Wang, and Mark Johnson. 2026.
\newblock {CORAL}: Cross-task correctness probes for robust LLM calibration.
\newblock In \emph{Proceedings of ACL}, pages 1823--1840.

\bibitem[Miao and Ungar(2026)]{miao2026orthogonality}
Qinyu Miao and Lyle Ungar. 2026.
\newblock Orthogonal correctness and confidence directions in language models: Implications for calibration.
\newblock In \emph{Proceedings of NAACL}, pages 2941--2958.

\bibitem[Nguyen et~al.(2021)]{nguyen2021wide}
Thao Nguyen, Maithra Raghu, and Simon Kornblith. 2021.
\newblock Do wide neural networks learn the same things? {C}onvergence over the training trajectory.
\newblock In \emph{Proceedings of ICLR}.

\bibitem[Pires et~al.(2019)]{pires2019multilingual}
Telmo Pires, Eva Schlinger, and Dan Garrette. 2019.
\newblock How multilingual is multilingual {BERT}?
\newblock In \emph{Proceedings of ACL}, pages 4996--5001.

\bibitem[Raghu et~al.(2017)]{raghu2017svcca}
Maithra Raghu, Justin Gilmer, Jason Yosinski, and Jascha Sohl-Dickstein. 2017.
\newblock {SVCCA}: Singular vector canonical correlation analysis for deep learning dynamics and interpretability.
\newblock In \emph{Proceedings of NeurIPS}, pages 6076--6085.

\bibitem[Voita and Titov(2020)]{voita2020mdl}
Elena Voita and Ivan Titov. 2020.
\newblock Information-theoretic probing with minimum description length.
\newblock In \emph{Proceedings of EMNLP}, pages 183--196.

\bibitem[Xiao et~al.(2025)]{xiao2025knowledgeboundary}
Shiyuan Xiao, Jiaxin Zhang, and Yun Chen. 2025.
\newblock Cross-lingual transfer of knowledge boundary probes in multilingual language models.
\newblock In \emph{Proceedings of ACL}, pages 5312--5328.

\bibitem[Zhou et~al.(2024)]{zhou2024lace}
Wenxuan Zhou, Sheng Zhang, Hoifung Poon, and Muhao Chen. 2024.
\newblock {LACE}: Layer-wise adaptive calibration ensembles for multilingual LLMs.
\newblock In \emph{Proceedings of EMNLP}, pages 7823--7838.

\end{thebibliography}

%====================================================================
\appendix

\section{INT8 Validation}
\label{app:int8}

We validate INT8 quantization against FP16 on Llama~4 Scout at 5 representative layers using a 5\% data subset. The maximum AUROC difference is 0.8\,pp, mean difference 0.3\,pp, with Pearson $r = 0.997$ between INT8 and FP16 probe scores. Rank agreement is 98\%. All differences fall within bootstrap confidence intervals, confirming INT8 does not distort correctness signals.

\section{Probe Complexity}
\label{app:probe_complexity}

MLP probes (2-layer) achieve 2.5\,pp higher English AUROC than logistic regression (0.835 vs.\ 0.81 on Qwen3) but suffer 6\,pp larger transfer gaps (0.17 vs.\ 0.11). The MLP exploits English-specific nonlinearities that do not generalize cross-lingually. 80\% of languages achieve better transfer with logistic regression, justifying our linear probe choice and reinforcing the geometric framing.

\section{Cross-Benchmark Generalization}
\label{app:benchmarks}

Results replicate on Belebele (EN AUROC: 0.78, ZS median: 0.68, $\rho$: $-0.48$) and MMLU-ProX (EN AUROC: 0.82, ZS median: 0.69, $\rho$: $-0.50$). Per-language AUROC rank correlations with GlobalMMLU are 0.82 (Belebele) and 0.88 (MMLU-ProX). Method rankings are consistent across all three benchmarks, confirming CaCLR's $\beta$ inertness is benchmark-independent.

\section{Control Task Selectivity}
\label{app:control}

Following \citet{hewitt2019control}, we train probes on randomly permuted labels. Control AUROC is 0.54--0.55 (near chance). Selectivity (probe AUROC minus control AUROC) is positive across all language--layer combinations: English median 26.5\,pp, high-resource 20\,pp, mid-resource 16\,pp, low-resource 10\,pp. 88\% of combinations exceed 10\,pp selectivity (95\% for high-resource, 70\% for low-resource). The weaker selectivity for low-resource languages is a genuine limitation: conclusions about sw/bn transfer should be interpreted with this caveat.

\end{document}
